\documentclass[twoside,openright,titlepage,fleqn,
	headinclude,11pt,a4paper,BCOR5mm,footinclude
	]{scrbook}
%\documentclass[a4paper,11pt]{book}
\usepackage[latin1]{inputenc}
\usepackage[T1]{fontenc}
\usepackage[italian]{babel}
\usepackage{graphicx}
\usepackage{subfiles}
\usepackage{listings}
%\usepackage{pgf}
\usepackage{xparse}
\usepackage{float}
\usepackage{color}
%\usepackage{xcolor}

%pacchetti del modello per la tesi
%--------------------------------------------------------------
\usepackage[square,numbers]{natbib} 
\usepackage[fleqn]{amsmath}  
\usepackage{dia-classicthesis-ldpkg} 
% Options for classicthesis.sty:
% tocaligned eulerchapternumbers drafting linedheaders 
% listsseparated subfig nochapters beramono eulermath parts 
% minionpro pdfspacing
\usepackage[eulerchapternumbers,subfig,beramono,eulermath,
	parts]{classicthesis}
%--------------------------------------------------------------
\usepackage{tikz}
\usetikzlibrary{shapes, chains, scopes, shadows, positioning, arrows}

%comandi per modello
%--------------------------------------------------------------
\newcommand{\myTitle}{Analisi e sviluppo di un crawler per la creazione di una base di conoscenza semantica del personale universitario\xspace}
% use the right myDegree option
\newcommand{\myDegree}{Corso di Laurea in Informatica\xspace}
%\newcommand{\myDegree}{
	%Corso di Laurea Specialistica in Scienze e Tecnologie 
	%dell'Informazione\xspace}
\newcommand{\myName}{Stefano Martina\xspace}
\newcommand{\myMail}{stefano.martina@gmail.com}
\newcommand{\myProf}{Elena Barcucci\xspace}
\newcommand{\myOtherProf}{Nome Cognome\xspace}
\newcommand{\mySupervisor}{Paolo Nesi\xspace}
\newcommand{\myFaculty}{
	Facolt\`a di Scienze Matematiche, Fisiche e Naturali\xspace}
\newcommand{\myDepartment}{
	Dipartimento di Sistemi e Informatica\xspace}
\newcommand{\myUni}{\protect{
	Universit\`a degli Studi di Firenze}\xspace}
\newcommand{\myLocation}{Firenze\xspace}
\newcommand{\myTime}{Anno Accademico 2011-2012\xspace}
\newcommand{\myVersion}{Version 0.1\xspace}
\newlength{\abcd} % for ab..z string length calculation
% how all the floats will be aligned
\newcommand{\myfloatalign}{\centering} 
\setlength{\extrarowheight}{3pt} % increase table row height
\captionsetup{format=hang,font=small}
%--------------------------------------------------------------
% Layout setting
%--------------------------------------------------------------
\usepackage{geometry}
\geometry{
	a4paper,
	ignoremp,
	bindingoffset = 1cm, 
	textwidth     = 13.5cm,
	textheight    = 21.5cm,
	lmargin       = 3.5cm, % left margin
	tmargin       = 4cm    % top margin 
}
%--------------------------------------------------------------

%comando per asterismo
\newcommand{\asterismo}{{\footnotesize \smash{% 
    \raisebox{-.2ex}{% 
      \setlength{\tabcolsep}{0.5pt}%
      \begin{tabular}{@{}cc@{}}% 
        \multicolumn2c*\\[-1.5ex] *&*% 
      \end{tabular}}}}}

%comando per impostazioni float di default
\makeatletter
\def\fps@figure{!htbp}
\def\fps@table{!htbp}
\def\fps@code{!htbp}
\makeatother

%cambia comportamento delle description
\renewcommand{\descriptionlabel}[1]{\hspace{2em}\hspace{\labelsep}\textbf{#1}}

%stili per disegnare con tikz
\colorlet{arancio}{yellow!20!red!20}
\colorlet{viola}{blue!10!red!5}

\tikzstyle{disp}=[rectangle, rounded corners=10pt, align=center, drop shadow, text width=2cm, fill=blue!20, draw=black!50]
\tikzstyle{db}=[cylinder, align=center, aspect=0.2, drop shadow, text width=2cm, shape border rotate=90, fill=yellow!20, draw=black!50]
\tikzstyle{rdf}=[cylinder, align=center, aspect=0.2, drop shadow, text width=2cm, shape border rotate=90, fill=yellow!20!red!20, draw=black!50]
\tikzstyle{web}=[rectangle, rounded corners=2pt, align=center, drop shadow, text width=2cm, fill=green!20, draw=black!50]
\tikzstyle{weblist}=[rectangle, rounded corners=2pt, align=center, double copy shadow, text width=2cm, fill=green!20, draw=black!50]
\tikzstyle{file}=[chamfered rectangle, align=center, drop shadow, text width=2cm, fill=black!20, draw=black!50]
\tikzstyle{filelist}=[chamfered rectangle, align=center, double copy shadow, text width=2cm, fill=black!20, draw=black!50]
\tikzstyle{gaz}=[tape, align=center, drop shadow, text width=2cm, shape border rotate=90, fill=blue!20, draw=black!50]
\tikzstyle{gate}=[rectangle, rounded corners=2pt, align=center, drop shadow, text width=2cm, fill=red!20, draw=black!50]
\tikzstyle{app}=[ellipse, align=center, drop shadow, text width=2cm, fill=red!20, draw=black!50]
\tikzstyle{esp}=[cloud, aspect=2, align=center, drop shadow, text width=2cm, fill=blue!10, draw=black!50]
\tikzstyle{rdfc}=[rectangle, rounded corners=5pt, align=center, drop shadow, text width=2cm, fill=green!20!red!20, draw=black!50]
\tikzstyle{rdfclist}=[rectangle, rounded corners=5pt, align=center, double copy shadow, text width=2cm, fill=green!20!red!20, draw=black!50]
\tikzstyle{webapp}=[ellipse, align=center, drop shadow, text width=2cm, fill=green!20, draw=black!50]
\tikzstyle{app2}=[rectangle, rounded corners=2pt, align=center, drop shadow, text width=2cm, fill=red!20, draw=black!50]
\tikzstyle{app3}=[rectangle, rounded corners=2pt, align=center, drop shadow, text width=2cm, fill=green!20, draw=black!50]
\tikzstyle{app4}=[rectangle, rounded corners=2pt, align=center, drop shadow, text width=2cm, fill=arancio, draw=black!50]
\tikzstyle{tab}=[rounded rectangle, align=center, drop shadow, text width=2cm, fill=blue!20, draw=black!50]
\tikzstyle{man}=[circle, align=center, drop shadow, text width=2cm, fill=red!20, draw=black!50]
\tikzstyle{label}=[rectangle, align=center, text width=2cm, rounded corners=2pt, fill=yellow!20]
\tikzstyle{round}=[circle, align=center, drop shadow, text width=0cm, fill=blue!20, draw=black!50]
\tikzstyle{choice}=[diamond, align=center, drop shadow, text width=2cm, fill=blue!20, draw=black!50]
\tikzstyle{legenda}=[rectangle, rounded corners=2pt, align=center, drop shadow, fill=viola, draw=black!50]
\tikzstyle{legendait}=[text width=10pt, text depth=10pt]
\tikzstyle{legendalab}=[]

\tikzstyle{freccia}=[->, very thick, >=stealth', draw=black!80]
\tikzstyle{tfreccia}=[->, dash pattern=on 3pt off2pt, very thick, >=stealth', draw=black!80]
\tikzstyle{dfreccia}=[<->, very thick, >=stealth', draw=black!80]
\tikzstyle{dtfreccia}=[<->, dash pattern=on 3pt off2pt, very thick, >=stealth', draw=black!80]
\tikzstyle{nfreccia}=[ very thick, >=stealth', draw=black!80]
\tikzstyle{tnfreccia}=[dash pattern=on 3pt off2pt, very thick, >=stealth', draw=black!80]

\tikzstyle{point}=[coordinate]
\tikzstyle{vuoto}=[text width=0pt]

\tikzstyle{forshadow}=[rectangle, rounded corners=2pt, drop shadow, inner sep = 0pt]
\tikzstyle{sqltable}=[rectangle split, rounded corners=2pt, inner sep = 2pt, rectangle split part align={center, left}, rectangle split parts=#1, draw, rectangle split part fill={green!20, yellow!20}, draw=black!50]

\colorlet{umlTit}{green!20}
\colorlet{umlAtt}{yellow!20!red!20}
\colorlet{umlMet}{yellow!20}
\tikzstyle{umltable}=[rectangle split, rounded corners=2pt, inner sep = 2pt, rectangle split part align={center, left}, rectangle split parts=#1, draw, rectangle split part fill={umlTit, umlMet}, draw=black!50]


\tikzstyle{coda}=[rectangle split, rectangle split part align={center}, rectangle split parts=#1, draw, fill=yellow!20, draw=black!50, drop shadow]

\tikzstyle{uri}=[ellipse, drop shadow, draw=black!50, fill=green!20, very thick]
\tikzstyle{letterale}=[rectangle, drop shadow, rounded corners=2pt, draw=black!50, fill=red!20, very thick]
\tikzstyle{arcol}=[->, >=stealth, draw=black!80, very thick, bend right]
\tikzstyle{labell}=[rectangle, rounded corners=2pt, fill=yellow!20, swap]
\tikzstyle{arcor}=[->, >=stealth, draw=black!80, very thick, bend left]
\tikzstyle{labelr}=[rectangle, rounded corners=2pt, fill=yellow!20]

\tikzstyle{uriEsistente}=[ellipse, drop shadow, draw=black!50, fill=blue!20, very thick]
\tikzstyle{uriChoice}=[rectangle, rounded corners=2pt, drop shadow, draw=black!50, fill=yellow!20, very thick, align=left]

\tikzstyle{labelg}=[rectangle, rounded corners=2pt, sloped, above, fill=yellow!20]

\tikzstyle{tempLab}=[rectangle, rounded corners=2pt, text width = 1.5cm, text height = 0.5cm, text depth = 0.25cm, align=center, drop shadow, fill=viola, draw=black!50]
\tikzstyle{tempLab2}=[rectangle, rounded corners=2pt, text width = 2cm, text height = 0.5cm, text depth = 0.25cm, align=center, drop shadow, fill=viola, draw=black!50]
\tikzstyle{tempLab3}=[rectangle, rounded corners=2pt, text width = 2cm, text height = 0.5cm, text depth = 0.75cm, align=center, drop shadow, fill=viola, draw=black!50]
\tikzstyle{tempHid3}=[rectangle, rounded corners=2pt, text width = 2cm, text height = 0.5cm, text depth = 0.75cm ]
\tikzstyle{tempLine}=[very thick, draw=black!80]
\tikzstyle{tempTLine}=[dash pattern=on 3pt off2pt, very thick, draw=black!80]
\tikzstyle{tempLineProc}=[very thick, draw=red!80]


%ambiente per astrarre il disegno di un grafo di triple
%senza stella usa l'unita di default di 8cm e i primi due parametri
%opzionali come caption e label.
%Con la stella usa il primo parametro come unita e i seguenti due come
%caption e label
\NewDocumentEnvironment{tgraph}{sooo}{
  \begin{figure}
    \begin{center}
      \IfBooleanTF #1{
        \IfValueTF{#2}{
          \begin{tikzpicture}[x=#2, y=#2, auto]
        }{
          \begin{tikzpicture}[x=8cm, y=8cm, auto]
        }
      }{
        \begin{tikzpicture}[x=8cm, y=8cm, auto]
      }
}{
        \end{tikzpicture}
   \end{center}
   \IfBooleanTF #1{
     \IfValueTF{#3}{\caption{#3}}{}
     \IfValueTF{#4}{\label{#4}}{}
   }{
     \IfValueTF{#2}{\caption{#2}}{}
     \IfValueTF{#3}{\label{#3}}{}
   }
 \end{figure}
}
\newcommand{\tgraphlet}[3]{
  \node (#1) at (#2) [letterale] {#3};
}
\newcommand{\tgraphuri}[3]{
  \node (#1) at (#2) [uri] {#3};
}
\newcommand{\tgrapharcl}[3]{
  \draw [arcol] (#1) to node[labell]{#3} (#2);
}
\newcommand{\tgrapharcr}[3]{
  \draw [arcor] (#1) to node[labelr]{#3} (#2);
}

%comando unico per tutti i codici
\newcommand{\unicocode}[1]{\texttt{\bfseries #1}}
\newcommand{\unicocodet}[1]{\texttt{\bfseries #1}}

%definisce un nuovo ambiente per un blocco java e un nuovo comando per
%un codice java in linea
\lstnewenvironment{javablock}
{
  \lstset{language=Java, basicstyle=\small\ttfamily\bfseries,
    breaklines=true, breakatwhitespace=false, frameround=fttt,
    frame=trBL, backgroundcolor=\color{yellow!20}}
}{}

\newcommand{\java}[1]{\unicocode{#1}}
%  ``\lstinline[language=Java,basicstyle=\small\ttfamily\bfseries,breaklines=true]|#1|''

\newcommand{\javat}[1]{\unicocode{#1}}
%  ``\lstinline[language=Java,basicstyle=\small\ttfamily\bfseries,breaklines=true]|#1|''

\newcommand{\javai}[1]{\unicocodet{#1}}
%  \lstinline[language=Java,basicstyle=\small\ttfamily\bfseries,breaklines=true]|#1|


%definisce un nuovo ambiente per un blocco xml e un nuovo comando per
%un codice xml in linea
\lstnewenvironment{xmlblock}
{
  \lstset{language=XML, basicstyle=\small\ttfamily\bfseries,
    breaklines=true, breakatwhitespace=false, frameround=fttt,
    frame=trBL, backgroundcolor=\color{yellow!20}}
}{}

\newcommand{\xml}[1]{\unicocode{#1}}
%  ``\lstinline[language=XML,basicstyle=\small\ttfamily\bfseries,breaklines=true]|#1|''


\newcommand{\xmli}[1]{\unicocodet{#1}}
%  \lstinline[language=XML,basicstyle=\small\ttfamily\bfseries,breaklines=true]|#1|


%definisce un nuovo ambiente per un blocco abnf e un nuovo comando per
%un codice abnf in linea
\lstnewenvironment{abnfblock}
{
  \lstset{basicstyle=\small\ttfamily\bfseries, commentstyle=\color{Black},
    breaklines=true, breakatwhitespace=false, frameround=fttt,
    frame=trBL, backgroundcolor=\color{yellow!20}}
}{}

\newcommand{\abnf}[1]{\unicocode{#1}}
%  ``\lstinline[basicstyle=\small\ttfamily\bfseries,breaklines=true]|#1|''


%definisce un nuovo ambiente per un blocco ntriples e un nuovo comando per
%un codice ntriples in linea
\lstnewenvironment{ntriplesblock}
{
  \lstset{language=XML, basicstyle=\small\ttfamily\bfseries,
    breaklines=true, breakatwhitespace=false, frameround=fttt,
    frame=trBL, backgroundcolor=\color{yellow!20}}
}{}

\newcommand{\ntriples}[1]{\unicocode{#1}}
%  ``\lstinline[language=XML,basicstyle=\small\ttfamily\bfseries,breaklines=true]|#1|''


%ambiente per inserire una singola tripla
\newcommand{\singletriple}[3]{
  \begin{center}
    \begin{tikzpicture}
        \node [rectangle, rounded corners=2pt, inner sep = 0.3cm, drop shadow, draw=black!50, fill=yellow!20] {
          \begin{tabular}{rcl}
            #1 & #2 & #3 \\
          \end{tabular}
        };
    \end{tikzpicture}
  \end{center}
}

%ambiente per disegnare un insieme di triple
\NewDocumentEnvironment{triples}{oo}{
  \begin{table}
    \begin{center}
      \begin{tabular}{|c|c|c|}
        \hline
        {\bfseries Soggetto} & {\bfseries Predicato} & {\bfseries Oggetto} \\
        \hline
}{
  \triple{}{}{}
  \hline
      \end{tabular}
    \end{center}
    \IfValueTF{#1}{\caption{#1}}{}
    \IfValueTF{#2}{\label{#2}}{}
  \end{table}
}

\newcommand{\triple}[3]{
  #1 & #2 & #3 \\
}

%comando per inserire una immagine
\NewDocumentCommand{\image}{ommm}{
  \begin{figure}
    \begin{center}
      \begin{tikzpicture}
        \node [rectangle, rounded corners=2pt, inner sep = 0.3cm, drop shadow, draw=black!50, fill=yellow!20] {
          \IfNoValueTF{#1}
                      {\includegraphics[width=13cm]{#2}}
                      {\includegraphics[width=#1]{#2}}
        };
      \end{tikzpicture}
    \end{center}
    \caption{#3}
    \label{#4}
  \end{figure}
}

\NewDocumentCommand{\imager}{ommm}{
  \begin{figure}
    \begin{center}
      \begin{tikzpicture}
        \node [rectangle, rounded corners=2pt, inner sep = 0.3cm, drop shadow, draw=black!50, fill=yellow!20, rotate = 90] {
          \IfNoValueTF{#1}
                      {\includegraphics[height=13cm]{#2}}
                      {\includegraphics[height=#1]{#2}}
        };
      \end{tikzpicture}
    \end{center}
    \caption{#3}
    \label{#4}
  \end{figure}
}

%comando per inserire un link
\newcommand{\link}[1]{\unicocode{#1}}


%comando per inserire un file
\newcommand{\file}[1]{\unicocode{#1}}

%comando per inserire un nome gate
\newcommand{\gate}[1]{\unicocode{#1}}

%comando per inserire dell'sql
\newcommand{\sql}[1]{\unicocode{#1}}
%  ``\lstinline[language=SQL,basicstyle=\small\ttfamily\bfseries,breaklines=true]|#1|''


\newcommand{\sqlt}[1]{\unicocode{#1}}
%  ``\lstinline[language=SQL,basicstyle=\small\ttfamily\bfseries,breaklines=true]|#1|''


\newcommand{\sqli}[1]{\unicocodet{#1}}
%  \lstinline[language=SQL,basicstyle=\small\ttfamily\bfseries,breaklines=true]|#1|


%definisce un nuovo ambiente float per il codice
\newfloat{code}{!htbp}{}
\floatname{code}{Codice}

%comando per dimensioni testo
\newcommand{\dimg}{\tiny}
\newcommand{\codg}[1]{\dimg \unicocodet{#1}}
%  \lstinline[basicstyle=\dimg\ttfamily\bfseries,breaklines=true]|#1|


\newcommand{\dims}{\scriptsize}
\newcommand{\cods}[1]{\dims\unicocodet{#1}}
%  \lstinline[basicstyle=\dims\ttfamily\bfseries,breaklines=true]|#1|


%\title{
%  \includegraphics[width=5cm]{img/logo-unifi.eps}\\[1cm]
%  \textbf{Analisi e sviluppo di un crawler per la creazione di una base di conoscenza semantica del personale universitario}
 %\\
  %parte del progetto \emph{OSIM} del \emph{Distributed Systems and Internet Technology Laboratory}
%}
%\author{\textbf{Stefano Martina}}
\begin{document}
%roba del modello
%--------------------------------------------------------------
\frenchspacing
\raggedbottom
\pagenumbering{Roman}
\pagestyle{plain}
%--------------------------------------------------------------
% Frontmatter
%--------------------------------------------------------------
%--------------------------------------------------------------
% titlepage.tex (use thesis.tex as main file)
%--------------------------------------------------------------
\begin{titlepage}
	\begin{center}
   	\large
      \hfill
      \vfill
      \begingroup
			\spacedallcaps{\myUni} \\ 
			\myFaculty \\
			\myDegree \\ 
			\vspace{0.5cm}
         \includegraphics[scale=.065]{logo/unifi}\\
         \vspace{0.5cm}    
         Tesi di Laurea    
      \endgroup 
      \vfill 
      \begingroup
      	\color{Maroon}\spacedallcaps{\myTitle} \\ \bigskip
      \endgroup
      \spacedlowsmallcaps{\myName}
      \vfill  
      Relatore: {\itshape\myProf}
      \vfill
      Corelatore: {\itshape\mySupervisor}
      \vfill
      \myTime
      \vfill                      
	\end{center}        
\end{titlepage}   
%--------------------------------------------------------------
% back titlepage
%--------------------------------------------------------------
   \newpage
	\thispagestyle{empty}
	\hfill
	\vfill
	\noindent\myName(\texttt{\myMail}): 
	\textit{\myTitle,} 
	\myDegree, \textcopyright\ \myTime
%--------------------------------------------------------------
% back titlepage end
%--------------------------------------------------------------

\cleardoublepage%*******************************************************
% Abstract
%*******************************************************
%\renewcommand{\abstractname}{Abstract}
%\pdfbookmark[1]{Abstract}{Abstract}
\begingroup
\let\clearpage\relax
\let\cleardoublepage\relax
\let\cleardoublepage\relax

\chapter*{Sommario}
Il presente testo \`e il risultato del lavoro di \emph{stage} svolto
presso il laboratorio \emph{DISIT} del Dipartimento di Sistemi e
Informatica dell'Universit\`a degli Studi di Firenze.

Il lavoro effettuato si inserisce nel progetto \emph{OSIM} che ha
l'obbiettivo di creare una base di conoscenza semantica riguardante il
personale universitario. In particolare \`e
stata sviluppata la parte dell'applicazione che riguarda il \emph{data
mining} delle informazioni riguardanti le persone, i corsi associati a
queste, e le competenze di ogni persona.

Il progetto fa uso delle
tecnologie \emph{RDF}, \emph{RDFS} e \emph{OWL} per l'organizzazione
della base di conoscenza semantica. Per quanto riguarda
l'applicazione, questa \`e sviluppata in \emph{Java} facendo uso delle
\emph{Servlet} e delle \emph{JSP} per fornire l'interfaccia \emph{web}
con l'utente.

Il testo \`e diviso in due parti: una parte consiste in una
sostanziosa lettura sullo stato dell'arte delle tecnologie usate,
dalle basi dell'\emph{XML} e delle tecnologie correlate sino ad una
visione eterogenea sulle tecnologie del web semantico;
l'altra parte concerne la descrizione del lavoro effettuato e un'analisi
dettagliata dell'applicazione nelle sue parti.

\endgroup			

\vfill

\cleardoublepage%*******************************************************
% Acknowledgments
%*******************************************************
%\pdfbookmark[1]{Acknowledgments}{acknowledgments}

\begin{flushright}{\slshape
    We can only see a short distance ahead, \\
    but we can see plenty there that needs to be done.} \\ \medskip
    --- Alan Turing
\end{flushright}



\bigskip

\begingroup
\let\clearpage\relax
\let\cleardoublepage\relax
\let\cleardoublepage\relax
\chapter*{Ringraziamenti}
Desidero ringraziare calorosamente la Prof. Elena Barcucci e il Prof. Paolo Nesi per
l'opportunit\`a offertami da questo lavoro e per l'aiuto datomi allo
sviluppo dell'applicazione e del presente testo. Ringrazio
particolarmente Ing. Alice Fuzier Cayla per il grande aiuto sulla
parte introduttiva a \emph{GATE} e Ph.D. Andrea Bellandi e Ph.D. Ing. Gianni Pantaleo
per il supporto dato allo sviluppo di alcune parti del programma e per il lavoro di revisione. 

Ringrazio inoltre Ph.D. Ing. Pierfrancesco Bellini e Ing. Nadia Rauch per l'aiuto datomi
all'inizio dello stage e Ph.D Ing. Michela Paolucci per il materiale
fornitomi. 
Vorrei ringraziare anche tutto il personale del laboratorio
\emph{DISIT} per l'ambiente dinamico, familiare e stimolante nel quale
ho avuto il piacere di lavorare in questi mesi.

Desidero infine ringraziare con grande affetto i miei genitori per
l'insostituibile supporto ricevuto in questi anni, e Federica per
aver reso questo testo pi\`u comprensibile e per la sua
pazienza nel sopportare la concorrenza del computer alla nostra
relazione\dots


\endgroup




\cleardoublepage\tableofcontents
%--------------------------------------------------------------
% Mainmatter
%--------------------------------------------------------------
\clearpage
\cleardoublepage
%\newpage
%\thispagestyle{empty}
%\mbox{}
%\newpage
%\thispagestyle{empty}
%\mbox{}
%--------------------------------------------------------------
\newpage
\pagenumbering{arabic}
\pagestyle{scrheadings}
\sloppy
\part{Stato dell'arte}
%\subfile{logica}
\subfile{xml}
\subfile{web_semantico}
\subfile{gate}
\part{Analisi e lavoro effettuato}
%\subfile{progetto_osim}
\subfile{analisi_crawler}
\subfile{analisi_applicazione}
\subfile{sviluppi_futuri}
\clearpage
\addcontentsline{toc}{chapter}{Bibliografia}
\begin{thebibliography}{9}

\bibitem{semantic-primer}Antoniou Grigorios, Harmelen Frank.
  \emph{A semantic web primer}.
  Cambridge MIT press, 2004.

\bibitem{working-ontologist}Dean Allemang, Jim Hendler.
  \emph{SEMANTIC WEB for the WORKING ONTOLOGIST}.
  Morgan Kaufmann, 2008.

\bibitem{xml-recomendation}W3C.
  \emph{Raccomandazioni del w3c relative a XML e alle tecnologie relative}.\\
  ``http://www.w3.org/TR/xml/'',\\
  ``http://www.w3.org/TR/xml-names/'',\\
  ``http://www.w3.org/TR/xmlschema-0/'',\\
  ``http://www.w3.org/TR/xmlschema-1/'',\\
  ``http://www.w3.org/TR/xmlschema-2/'',\\
  ``http://www.w3.org/TR/xslt/'',\\
  ``http://www.w3.org/TR/xpath/'',\\
  ``http://www.w3.org/TR/xsl/'',\\
  ``http://www.w3.org/TR/xml-infoset''.

\bibitem{rdf-recomendation}W3C.
  \emph{Raccomandazioni del w3c relative a RDF}.\\
  ``http://www.w3.org/TR/rdf-syntax-grammar/'',\\
  ``http://www.w3.org/TR/rdf-primer/'',\\
  ``http://www.w3.org/TR/rdf-concepts/'',\\
  ``http://www.w3.org/TR/rdf-syntax-grammar/'',\\
  ``http://www.w3.org/TR/rdf-mt/'',\\
  ``http://www.w3.org/TR/rdf-schema/'',\\
  ``http://www.w3.org/TR/rdf-testcases/''.

\bibitem{owl}W3C.
  \emph{Raccomandazioni del w3c relative a OWL}.\\
  ``http://www.w3.org/TR/owl-guide/'',\\
  ``http://www.w3.org/TR/owl-features/''.

\bibitem{skos}W3C.
  \emph{Raccomandazioni del w3c relative a SKOS}.\\
  ``http://www.w3.org/TR/skos-primer/'',\\
  ``http://www.w3.org/TR/skos-reference/'',\\
  ``http://www.w3.org/2004/02/skos/''.

\bibitem{foaf}D. Brickley, L. Miller.
  \emph{FOAF Vocabulary Specification}.\\
  ``http://xmlns.com/foaf/spec/''.

\bibitem{std66}T. Berners Lee, et al.
  \emph{STD 66, RFC 3986, Uniform Resource Identifier (URI): Generic
    Syntax}.\\
  ``http://tools.ietf.org/html/std66'', January 2005.

\bibitem{std68}D. Crocker, P. Overell.
  \emph{STD 68, RFC 5234, Augmented BNF for Syntax Specifications: ABNF}.\\
  ``http://tools.ietf.org/html/std68'', January 2008.

\bibitem{jqueryuilayout} \emph{jQuery UI Layout}.\\
  ``http://layout.jquery-dev.net/''

\bibitem{Cunningham2011a}Hamish Cunningham and Diana Maynard and Kalina Bontcheva and Valentin Tablan and Niraj Aswani and Ian Roberts and Genevieve Gorrell and Adam Funk and Angus Roberts and Danica Damljanovic and Thomas Heitz and Mark A. Greenwood and Horacio Saggion and Johann Petrak and Yaoyong Li and Wim Peters.
 \emph{Text Processing with GATE (Version 6)}. 2011.\\
 ``http://tinyurl.com/gatebook''

\bibitem{Appelt:1998:CPS:1119089.1119095}Appelt, Douglas E. and Onyshkevych, Boyan.
 \emph{The common pattern specification language}, 1998.

\bibitem{Cunningham2002}Hamish Cunningham and Diana Maynard and Kalina Bontcheva and Valentin Tablan.
\emph{GATE: A Framework and Graphical Development Environment for Robust NLP Tools and Applications}. 2002.

\end{thebibliography}

\end{document}
